% GENERATED FILE. Edit canonical lessons, then run npm run generate:books.
% canonical-source-hash: fnv1a64:83a14a50570e3508
% canonical-lessons: TA-C262-avvappotu, TA-C262-varantorum, TA-C262-tamatamaka, TA-C262-kataiciyaka, TA-C262-ippotellam, TA-R262-first-pass-words-from-chapters-254-257, TA-R262-second-pass-words-from-chapters-258-262

\chapter{Now and then, Every week, Late, Finally, These days}
\label{ch:ta-now-and-then-every-week-late-finally-these-days}
\hlchaptermodality{262}

\begin{chapteropening}
\textbf{By the end of this chapter:} \emph{I can say now and then, every week, late, finally and these days.}

Complete the last lesson of chapter 262: I can say now and then, every week, late, finally and these days.
\end{chapteropening}

\section[\texorpdfstring{\ta{அவ்வப்போது} (avvappōtu)}{அவ்வப்போது (avvappōtu)}]{\ta{அவ்வப்போது} (avvappōtu) — now and then, occasionally}

\label{lesson:TA-C262-avvappotu}



\begin{quote}
Before the new one: say the Tamil for an itch, then the Tamil for dizziness.
\end{quote}

\subsection*{You'll want to know: \ta{அவ்வப்போது}}
\textbf{\ta{அவ்வப்போது}} — \emph{avvappōtu} — \textquotedblleft{}now and then, occasionally\textquotedblright{}.

\begin{grammarlens}[title={What you've built}]
One more word for this chapter.
\end{grammarlens}

\subsection*{Your turn}
\begin{itemize}
\raggedright
  \item \emph{Say it:} \emph{avvappōtu}
  \item \emph{Say it:} \emph{avvappōtu}, once more
  \item \emph{Say it:} say \emph{avvappōtu}
  \item \emph{Recall:} say the Tamil for an itch, then the Tamil for dizziness, then say \emph{avvappōtu} again
\end{itemize}

\begin{tcolorbox}[breakable,colback=teal!4,colframe=teal!35!black,title={Before you move on}]
What does \ta{அவ்வப்போது} mean? (\textquotedblleft{}Now and then, occasionally\textquotedblright{}.) Say it once more.
\end{tcolorbox}
\section[\texorpdfstring{\ta{வாரந்தோறும்} (vārantōṟum)}{வாரந்தோறும் (vārantōṟum)}]{\ta{வாரந்தோறும்} (vārantōṟum) — every week}

\label{lesson:TA-C262-varantorum}



\begin{quote}
Before the new one: say the Tamil for dizziness, then the Tamil for now and then.
\end{quote}

\subsection*{You'll want to know: \ta{வாரந்தோறும்}}
\textbf{\ta{வாரந்தோறும்}} — \emph{vārantōṟum} — \textquotedblleft{}every week\textquotedblright{}.

\begin{grammarlens}[title={What you've built}]
One more word for this chapter.
\end{grammarlens}

\subsection*{Your turn}
\begin{itemize}
\raggedright
  \item \emph{Say it:} \emph{vārantōṟum}
  \item \emph{Say it:} \emph{vārantōṟum}, once more
  \item \emph{Say it:} say \emph{vārantōṟum}
  \item \emph{Recall:} say the Tamil for dizziness, then the Tamil for now and then, then say \emph{vārantōṟum} again
\end{itemize}

\begin{tcolorbox}[breakable,colback=teal!4,colframe=teal!35!black,title={Before you move on}]
What does \ta{வாரந்தோறும்} mean? (\textquotedblleft{}Every week\textquotedblright{}.) Say it once more.
\end{tcolorbox}
\section[\texorpdfstring{\ta{தாமதமாக} (tāmatamāka)}{தாமதமாக (tāmatamāka)}]{\ta{தாமதமாக} (tāmatamāka) — late}

\label{lesson:TA-C262-tamatamaka}



\begin{quote}
Before the new one: say the Tamil for now and then, then the Tamil for every week.
\end{quote}

\subsection*{You'll want to know: \ta{தாமதமாக}}
\textbf{\ta{தாமதமாக}} — \emph{tāmatamāka} — \textquotedblleft{}late\textquotedblright{}.

\begin{grammarlens}[title={What you've built}]
One more word for this chapter.
\end{grammarlens}

\subsection*{Your turn}
\begin{itemize}
\raggedright
  \item \emph{Say it:} \emph{tāmatamāka}
  \item \emph{Say it:} \emph{tāmatamāka}, once more
  \item \emph{Say it:} say \emph{tāmatamāka}
  \item \emph{Recall:} say the Tamil for now and then, then the Tamil for every week, then say \emph{tāmatamāka} again
\end{itemize}

\begin{tcolorbox}[breakable,colback=teal!4,colframe=teal!35!black,title={Before you move on}]
What does \ta{தாமதமாக} mean? (\textquotedblleft{}Late\textquotedblright{}.) Say it once more.
\end{tcolorbox}
\section[\texorpdfstring{\ta{கடைசியாக} (kaṭaiciyāka)}{கடைசியாக (kaṭaiciyāka)}]{\ta{கடைசியாக} (kaṭaiciyāka) — finally, at last}

\label{lesson:TA-C262-kataiciyaka}



\begin{quote}
Before the new one: say the Tamil for every week, then the Tamil for late.
\end{quote}

\subsection*{You'll want to know: \ta{கடைசியாக}}
\textbf{\ta{கடைசியாக}} — \emph{kaṭaiciyāka} — \textquotedblleft{}finally, at last\textquotedblright{}.

\begin{grammarlens}[title={What you've built}]
One more word for this chapter.
\end{grammarlens}

\subsection*{Your turn}
\begin{itemize}
\raggedright
  \item \emph{Say it:} \emph{kaṭaiciyāka}
  \item \emph{Say it:} \emph{kaṭaiciyāka}, once more
  \item \emph{Say it:} say \emph{kaṭaiciyāka}
  \item \emph{Recall:} say the Tamil for every week, then the Tamil for late, then say \emph{kaṭaiciyāka} again
\end{itemize}

\begin{tcolorbox}[breakable,colback=teal!4,colframe=teal!35!black,title={Before you move on}]
What does \ta{கடைசியாக} mean? (\textquotedblleft{}Finally, at last\textquotedblright{}.) Say it once more.
\end{tcolorbox}
\section[\texorpdfstring{\ta{இப்போதெல்லாம்} (ippōtellām)}{இப்போதெல்லாம் (ippōtellām)}]{\ta{இப்போதெல்லாம்} (ippōtellām) — these days, nowadays}

\label{lesson:TA-C262-ippotellam}



\begin{quote}
Before the new one: say the Tamil for late, then the Tamil for finally.
\end{quote}

\subsection*{You'll want to know: \ta{இப்போதெல்லாம்}}
\textbf{\ta{இப்போதெல்லாம்}} — \emph{ippōtellām} — \textquotedblleft{}these days, nowadays\textquotedblright{}.

\begin{grammarlens}[title={What you've built}]
One more word for this chapter.
\end{grammarlens}

\subsection*{Your turn}
\begin{itemize}
\raggedright
  \item \emph{Say it:} \emph{ippōtellām}
  \item \emph{Say it:} \emph{ippōtellām}, once more
  \item \emph{Say it:} say \emph{ippōtellām}
  \item \emph{Recall:} say the Tamil for late, then the Tamil for finally, then say \emph{ippōtellām} again
\end{itemize}

\begin{tcolorbox}[breakable,colback=teal!4,colframe=teal!35!black,title={Before you move on}]
What does \ta{இப்போதெல்லாம்} mean? (\textquotedblleft{}These days, nowadays\textquotedblright{}.) Say it once more.
\end{tcolorbox}
\section[(dialogue)]{Words from chapters 254-257 — a first pass}

\label{lesson:TA-R262-first-pass-words-from-chapters-254-257}



\begin{quote}
Say the words for finally and these days.
\end{quote}

\subsection*{Your turn}
Say these aloud:

\begin{itemize}
\raggedright
  \item to borrow, to pay, to choose, to enter, to exit
  \item to breathe, to recover, to swell, to wake, to brush one's teeth
  \item to accept, to follow, to miss, to erase, to print
  \item to hate, to be surprised, to feel shy, to be jealous, to kiss
\end{itemize}

\begin{tcolorbox}[breakable,colback=teal!4,colframe=teal!35!black,title={Before you move on}]
To borrow? (\textbf{\ta{கடன்} \ta{வாங்கு}}, \emph{kaṭaṉ vāṅku}.) To pay? (\textbf{\ta{செலுத்து}}, \emph{celuttu}.) To choose? (\textbf{\ta{தேர்ந்தெடு}}, \emph{tērnteṭu}.) To enter? (\textbf{\ta{நுழை}}, \emph{nuḻai}.) To exit? (\textbf{\ta{வெளியேறு}}, \emph{veḷiyēṟu}.) To breathe? (\textbf{\ta{சுவாசி}}, \emph{cuvāci}.) To recover? (\textbf{\ta{குணமாகு}}, \emph{kuṇamāku}.) To swell? (\textbf{\ta{வீங்கு}}, \emph{vīṅku}.) To wake? (\textbf{\ta{எழுப்பு}}, \emph{eḻuppu}.) To brush one's teeth? (\textbf{\ta{பல்} \ta{துலக்கு}}, \emph{pal tulakku}.) To accept? (\textbf{\ta{ஏற்றுக்கொள்}}, \emph{ēṟṟukkoḷ}.) To follow? (\textbf{\ta{பின்பற்று}}, \emph{piṉpaṟṟu}.) To miss? (\textbf{\ta{தவறவிடு}}, \emph{tavaṟaviṭu}.) To erase? (\textbf{\ta{அழி}}, \emph{aḻi}.) To print? (\textbf{\ta{அச்சிடு}}, \emph{acciṭu}.) To hate? (\textbf{\ta{வெறு}}, \emph{veṟu}.) To be surprised? (\textbf{\ta{ஆச்சரியப்படு}}, \emph{āccariyappaṭu}.) To feel shy? (\textbf{\ta{வெட்கப்படு}}, \emph{veṭkappaṭu}.) To be jealous? (\textbf{\ta{பொறாமைப்படு}}, \emph{poṟāmaippaṭu}.) To kiss? (\textbf{\ta{முத்தமிடு}}, \emph{muttamiṭu}.)
\end{tcolorbox}
\section[(dialogue)]{Words from chapters 258-262 — a second pass}

\label{lesson:TA-R262-second-pass-words-from-chapters-258-262}



\begin{quote}
Say the words for to melt and to pick.
\end{quote}

\subsection*{Your turn}
Say these aloud:

\begin{itemize}
\raggedright
  \item to melt, to pick, to bark, to comfort, to shake hands
  \item to take a photo, to act, to enjoy, to repair, to tidy up
  \item to increase, to spread, to sink, to float, to squeeze
  \item diabetes, exercise, yoga, an itch, dizziness
  \item now and then, every week, late, finally, these days
\end{itemize}

\begin{tcolorbox}[breakable,colback=teal!4,colframe=teal!35!black,title={Before you move on}]
To melt? (\textbf{\ta{உருகு}}, \emph{uruku}.) To pick? (\textbf{\ta{பறி}}, \emph{paṟi}.) To bark? (\textbf{\ta{குரை}}, \emph{kurai}.) To comfort? (\textbf{\ta{தேற்று}}, \emph{tēṟṟu}.) To shake hands? (\textbf{\ta{கை} \ta{குலுக்கு}}, \emph{kai kulukku}.) To take a photo? (\textbf{\ta{படம்} \ta{எடு}}, \emph{paṭam eṭu}.) To act? (\textbf{\ta{நடி}}, \emph{naṭi}.) To enjoy? (\textbf{\ta{ரசி}}, \emph{raci}.) To repair? (\textbf{\ta{பழுது} \ta{பார்}}, \emph{paḻutu pār}.) To tidy up? (\textbf{\ta{ஒழுங்குபடுத்து}}, \emph{oḻuṅkupaṭuttu}.) To increase? (\textbf{\ta{அதிகரி}}, \emph{atikari}.) To spread? (\textbf{\ta{பரவு}}, \emph{paravu}.) To sink? (\textbf{\ta{மூழ்கு}}, \emph{mūḻku}.) To float? (\textbf{\ta{மித}}, \emph{mita}.) To squeeze? (\textbf{\ta{பிழி}}, \emph{piḻi}.) Diabetes? (\textbf{\ta{சர்க்கரை} \ta{நோய்}}, \emph{carkkarai nōy}.) Exercise? (\textbf{\ta{உடற்பயிற்சி}}, \emph{uṭaṟpayiṟci}.) Yoga? (\textbf{\ta{யோகா}}, \emph{yōkā}.) An itch? (\textbf{\ta{அரிப்பு}}, \emph{arippu}.) Dizziness? (\textbf{\ta{தலைச்சுற்றல்}}, \emph{talaiccuṟṟal}.) Now and then? (\textbf{\ta{அவ்வப்போது}}, \emph{avvappōtu}.) Every week? (\textbf{\ta{வாரந்தோறும்}}, \emph{vārantōṟum}.) Late? (\textbf{\ta{தாமதமாக}}, \emph{tāmatamāka}.) Finally? (\textbf{\ta{கடைசியாக}}, \emph{kaṭaiciyāka}.) These days? (\textbf{\ta{இப்போதெல்லாம்}}, \emph{ippōtellām}.)
\end{tcolorbox}
